\documentclass[11pt]{article}

\usepackage[margin=1in]{geometry}
\usepackage{amsmath,amssymb,amsthm,mathtools}
\usepackage[colorlinks=true,linkcolor=blue,urlcolor=blue]{hyperref}
\usepackage{array,booktabs,tabularx}
\usepackage{enumitem}

\newtheorem{theorem}{Theorem}[section]
\newtheorem{proposition}[theorem]{Proposition}
\newtheorem{lemma}[theorem]{Lemma}
\newtheorem{corollary}[theorem]{Corollary}
\theoremstyle{definition}
\newtheorem{definition}[theorem]{Definition}
\newtheorem{conjecture}[theorem]{Conjecture}
\newtheorem{remark}[theorem]{Remark}

\newcommand{\R}{\mathbb R}
\newcommand{\N}{\mathbb N}
\newcommand{\rhop}{[\rho]_+}
\newcommand{\Fhat}{\widehat F}
\newcommand{\M}{\mathcal M}
\newcommand{\D}{\mathcal D}
\newcommand{\T}{\mathcal T}
\newcolumntype{Y}{>{\raggedright\arraybackslash}X}

\title{A Corrected Integrodifference Competition Model\\
in a Shifting Habitat}
\author{Xiang Huang \and Liang Kong}
\date{July 27, 2026}

\begin{document}

\maketitle

\begin{abstract}
We formulate a positivity-preserving integrodifference competition model in
a habitat moving at constant speed.  Signed environmental variation is
placed in the low-density growth factor, while density dependence remains
nonnegative.  We establish the local order structure, invariance of the
unit-square state space, preservation of the competitive product order, the
exact action of dispersal on exponential modes, an exponential comparison
principle, and the geometry of the favorable moving corridor.  We then give
precise conjectural formulations of the remaining extinction, exclusion, and
coexistence results.  An appendix audits the original proposal and records,
by category, which claims fail as stated and how they can be repaired.
\end{abstract}

\section{Introduction}

Integrodifference equations are natural models for populations with
discrete generations and nonlocal dispersal.  If favorable habitat moves at
speed \(c>0\), persistence requires the population to disperse and grow fast
enough to remain in the favorable region.  For two competing species, the
same problem also depends on invasion at the two semi-trivial equilibria and
on the order structure of the nonlinear update.

The purpose of this note is twofold.  First, we give a corrected model for
which nonnegative initial data generate nonnegative solutions and the
two-species update is genuinely competitive.  Second, we separate the
conclusions that follow from elementary comparison from the global
persistence statements that still require moving subsolutions and interior
squeezing.  The distinction is important: empty-habitat spreading speeds
control upper fronts, but they do not by themselves determine coexistence or
competitive exclusion.

\section{Model and assumptions}

For \(i=1,2\), let \(k_i\colon\R\to[0,\infty)\) be a continuous integrable
probability density,
\[
  \int_{\R}k_i(z)\,dz=1.
\]
Let \(\rho_i\colon\R\to(-1,\infty)\) be continuous and nondecreasing, with
finite limits
\[
  -1<\rho_i^-:=\lim_{s\to-\infty}\rho_i(s)<0
  <\rho_i^+:=\lim_{s\to+\infty}\rho_i(s).
\]
The species may have distinct environmental responses; the specialization
\(\rho_1=\rho_2\) is allowed.  Let \(\alpha_1,\alpha_2\geq0\).

For \(\rho>-1\) and \(\alpha,u,v\geq0\), define
\begin{equation}\label{eq:corrected-response}
  \Fhat(\rho,\alpha;u,v)
  :=
  \frac{(1+\rho)u}
       {1+[\rho]_+(u+\alpha v)},
  \qquad
  [\rho]_+:=\max\{\rho,0\}.
\end{equation}
The corrected two-species system is
\begin{align}
  u_{n+1}(x)
  &=
  \int_{\R}k_1(x-y)
  \Fhat\!\left(
    \rho_1(y-cn),\alpha_1;u_n(y),v_n(y)
  \right)\,dy,
  \label{eq:corrected-u}\\
  v_{n+1}(x)
  &=
  \int_{\R}k_2(x-y)
  \Fhat\!\left(
    \rho_2(y-cn),\alpha_2;v_n(y),u_n(y)
  \right)\,dy.
  \label{eq:corrected-v}
\end{align}
In favorable habitat, where \(\rho\geq0\), this is the usual
Beverton--Holt competition response
\[
  \Fhat(\rho,\alpha;u,v)
  =
  \frac{(1+\rho)u}{1+\rho(u+\alpha v)}.
\]
In unfavorable habitat, where \(-1<\rho\leq0\), it becomes
\[
  \Fhat(\rho,\alpha;u,v)=(1+\rho)u.
\]
Thus unfavorable habitat gives subcritical linear growth instead of
negative density dependence.

\section{Local dynamics and competitive order}

\begin{proposition}[Positivity, bounds, and local order]
\label{prop:local-order}
Suppose \(\rho>-1\) and \(\alpha,u,v\geq0\).  Then
\begin{enumerate}[label=(\roman*)]
  \item \(\Fhat(\rho,\alpha;u,v)\geq0\);
  \item \(\Fhat(\rho,\alpha;u,v)\leq(1+\rho)u\);
  \item if \(0\leq u\leq1\), then
        \(0\leq\Fhat(\rho,\alpha;u,v)\leq1\);
  \item \(\Fhat\) is nondecreasing in \(u\) and nonincreasing in \(v\).
\end{enumerate}
\end{proposition}

\begin{proof}
Put \(b=[\rho]_+\) and \(D=1+b(u+\alpha v)\).  Since
\(b,u,\alpha,v\geq0\), one has \(D\geq1\).  The numerator is nonnegative
because \(1+\rho>0\), which proves (i).  Dividing the numerator by
\(D\geq1\) proves (ii).

For (iii), suppose first that \(\rho\leq0\).  Then \(b=0\), and hence
\[
  \Fhat(\rho,\alpha;u,v)=(1+\rho)u\leq u\leq1.
\]
If \(\rho\geq0\), then
\[
  D-(1+\rho)u=(1-u)+\rho\alpha v\geq0,
\]
which again gives \(\Fhat\leq1\).

Finally, on the positive cone,
\[
  \frac{\partial\Fhat}{\partial u}
  =
  \frac{(1+\rho)(1+b\alpha v)}{D^2}\geq0,
  \qquad
  \frac{\partial\Fhat}{\partial v}
  =
  -\frac{(1+\rho)ub\alpha}{D^2}\leq0.
\]
This proves (iv).
\end{proof}

\begin{proposition}[Favorable coexistence equilibrium]
\label{prop:equilibrium}
Suppose \(0<\alpha_1,\alpha_2<1\), and define
\[
  u^*=\frac{1-\alpha_1}{1-\alpha_1\alpha_2},
  \qquad
  v^*=\frac{1-\alpha_2}{1-\alpha_1\alpha_2}.
\]
Then \(u^*,v^*>0\),
\[
  u^*+\alpha_1v^*=1,
  \qquad
  v^*+\alpha_2u^*=1,
\]
and for every \(\rho\geq0\),
\[
  \Fhat(\rho,\alpha_1;u^*,v^*)=u^*,
  \qquad
  \Fhat(\rho,\alpha_2;v^*,u^*)=v^*.
\]
\end{proposition}

\begin{proof}
Since \(0<\alpha_1\alpha_2<1\), all displayed denominators are positive.
Direct substitution gives the two nullcline identities.  If
\(u+\alpha v=1\) and \(\rho\geq0\), then
\[
  \Fhat(\rho,\alpha;u,v)
  =\frac{(1+\rho)u}{1+\rho}=u.
\]
Applying this identity to the two nullclines proves the result.
\end{proof}

The low-density multiplier of a rare species in favorable habitat, when its
competitor is at density one, is
\begin{equation}\label{eq:rare-multiplier}
  m(\rho,\alpha):=\frac{1+\rho}{1+\rho\alpha},
  \qquad \rho>0.
\end{equation}

\begin{proposition}[Weak and strong competition at a resident equilibrium]
\label{prop:invasion-multiplier}
Let \(\rho>0\).
\begin{enumerate}[label=(\roman*)]
  \item If \(0\leq\alpha<1\), then \(m(\rho,\alpha)>1\).
  \item If \(\alpha>1\), then \(m(\rho,\alpha)<1\).
\end{enumerate}
\end{proposition}

\begin{proof}
The denominator in \eqref{eq:rare-multiplier} is positive.  The first
inequality is equivalent to
\(\rho\alpha<\rho\), and the second is equivalent to
\(\rho<\rho\alpha\).
\end{proof}

For bounded continuous states, write
\[
  (u_1,v_1)\preceq_{\rm c}(u_2,v_2)
  \quad\Longleftrightarrow\quad
  u_1(x)\leq u_2(x)
  \ \text{and}\
  v_1(x)\geq v_2(x)
  \quad\text{for every }x\in\R.
\]
This is the competitive product order.  Let \(\T_n\) denote one update of
\eqref{eq:corrected-u}--\eqref{eq:corrected-v}.

\begin{theorem}[State-space invariance and competitive comparison]
\label{thm:state-space}
Assume the kernels and environments satisfy the hypotheses of Section~2.
Then \(\T_n\) maps bounded continuous
\([0,1]^2\)-valued states to bounded continuous
\([0,1]^2\)-valued states.  Moreover,
\[
  (u_1,v_1)\preceq_{\rm c}(u_2,v_2)
  \quad\Longrightarrow\quad
  \T_n(u_1,v_1)\preceq_{\rm c}\T_n(u_2,v_2).
\]
\end{theorem}

\begin{proof}
Proposition~\ref{prop:local-order} bounds each source term between zero and
one.  Integration against a nonnegative kernel of mass one preserves these
bounds.  Bounded continuity of the new profiles follows from the standard
continuity theorem for convolution of an \(L^1\) kernel with a bounded
continuous function.

For the order statement, increasing the focal profile and decreasing the
competitor profile increases the corresponding local source by
Proposition~\ref{prop:local-order}(iv).  Multiplication by a nonnegative
kernel and integration preserve the inequality.  Applying this once to the
\(u\)-equation and once, with the two order directions reversed, to the
\(v\)-equation proves the claim.
\end{proof}

\begin{corollary}[An absent species remains absent]
\label{cor:zero-species}
If \(v_0\equiv0\), then \(v_n\equiv0\) for every \(n\in\N\).  The analogous
statement holds for \(u\).
\end{corollary}

\begin{proof}
Equation~\eqref{eq:corrected-response} gives
\(\Fhat(\rho,\alpha;0,u)=0\).  Hence \(v_n\equiv0\) implies
\(v_{n+1}\equiv0\), and induction completes the proof.
\end{proof}

\section{Dispersal, exponential modes, and upper barriers}

For a kernel \(K\), define
\[
  (\D_Kf)(x):=\int_{\R}K(x-y)f(y)\,dy,
  \qquad
  \M_K(\mu):=\int_{\R}K(z)e^{\mu z}\,dz.
\]

\begin{lemma}[Exponential eigenfunction identity]
\label{lem:eigenfunction}
Whenever \(\M_K(\mu)\) is finite,
\[
  \D_K(e^{-\mu\,\cdot})(x)
  =
  \M_K(\mu)e^{-\mu x}.
\]
\end{lemma}

\begin{proof}
With \(z=x-y\),
\begin{align*}
  \int_{\R}K(x-y)e^{-\mu y}\,dy
  &=
  \int_{\R}K(z)e^{-\mu(x-z)}\,dz\\
  &=
  e^{-\mu x}\int_{\R}K(z)e^{\mu z}\,dz.
\end{align*}
\end{proof}

\begin{proposition}[Exact exponential orbit]
\label{prop:exponential-orbit}
Suppose \(\M_K(\mu)\) is finite.  Let \(Lf=R\D_Kf\), and let
\[
  W_n(x)=A\bigl(R\M_K(\mu)\bigr)^n e^{-\mu x}.
\]
Then \(W_{n+1}=LW_n\).  Along \(x=sn\),
\[
  W_n(sn)
  =
  A\left(R\M_K(\mu)e^{-\mu s}\right)^n.
\]
If \(R\M_K(\mu)>0\), \(\mu>0\), and
\[
  s>\frac{\log(R\M_K(\mu))}{\mu},
\]
then \(W_n(sn)\to0\).
\end{proposition}

\begin{proof}
The recurrence follows from Lemma~\ref{lem:eigenfunction}.  The moving-frame
formula is an algebraic rearrangement.  The speed inequality is equivalent
to
\[
  0<R\M_K(\mu)e^{-\mu s}<1,
\]
so the geometric factor tends to zero.
\end{proof}

\begin{theorem}[Exponential comparison]
\label{thm:barrier}
Assume \(K\geq0\) and \(R>0\).  Fix \(\mu>0\) with
\(0<\M_K(\mu)<\infty\), and suppose a nonnegative orbit \(w_n\) satisfies
\[
  w_{n+1}(x)\leq R\D_Kw_n(x).
\]
Suppose \(A\geq0\) and
\[
  w_0(x)\leq Ae^{-\mu x}
  \qquad (x\in\R).
\]
Then
\[
  w_n(x)\leq
  A\bigl(R\M_K(\mu)\bigr)^n e^{-\mu x}
  \qquad (n\in\N,\ x\in\R).
\]
Consequently, if
\[
  s>\frac{\log(R\M_K(\mu))}{\mu},
\]
then
\[
  \sup_{x\geq sn}w_n(x)\longrightarrow0.
\]
\end{theorem}

\begin{proof}
The dispersal operator preserves pointwise order because \(K\geq0\).
Induct on \(n\).  The initial inequality is the base case.  If the bound
holds at time \(n\), then
\[
  w_{n+1}
  \leq R\D_Kw_n
  \leq R\D_KW_n
  =W_{n+1}.
\]
For \(x\geq sn\), the exponential is decreasing, so
\[
  0\leq w_n(x)
  \leq
  A\left(R\M_K(\mu)e^{-\mu s}\right)^n.
\]
The right-hand side is a geometric sequence with ratio strictly between
zero and one.
\end{proof}

\begin{corollary}[Componentwise scalar domination]
\label{cor:scalar-domination}
Suppose \(\rho_i(s)\leq\rho_i^+\) for all \(s\), and set
\(R_i=1+\rho_i^+\).  Each component of
\eqref{eq:corrected-u}--\eqref{eq:corrected-v} satisfies
\[
  u_{n+1}\leq R_1\D_{k_1}u_n,
  \qquad
  v_{n+1}\leq R_2\D_{k_2}v_n.
\]
Thus Theorem~\ref{thm:barrier} applies separately to both species.
\end{corollary}

\begin{proof}
By Proposition~\ref{prop:local-order},
\[
  \Fhat(\rho_i(s),\alpha;u,v)
  \leq(1+\rho_i(s))u
  \leq(1+\rho_i^+)u.
\]
Multiplication by a nonnegative kernel and integration preserve the
inequality.
\end{proof}

\begin{remark}
Every bounded, compactly supported nonnegative initial profile lies below
\(Ae^{-\mu x}\) for a suitable \(A\).  Hence the preceding upper-front
estimate applies to the initial data normally used in spreading theory.
It is an upper-spreading result; persistence below the critical speed
requires a separate lower-solution argument.
\end{remark}

\section{The favorable moving corridor}

\begin{lemma}[Moving-coordinate geometry]
\label{lem:moving-coordinate}
Let the habitat move at speed \(c\).
\begin{enumerate}[label=(\roman*)]
  \item If an observer moves at speed \(s<c\), then
        \((s-c)n\to-\infty\).
  \item If \(s>c\), then \((s-c)n\to+\infty\).
  \item If \(x_n\geq(c+\varepsilon)n\) for some \(\varepsilon>0\), then
        \(x_n-cn\to+\infty\).
\end{enumerate}
\end{lemma}

\begin{proof}
The first two statements follow from the sign of \(s-c\).  For the third,
\[
  x_n-cn\geq\varepsilon n\longrightarrow+\infty.
\]
\end{proof}

\begin{corollary}[Uniformly favorable corridor geometry]
\label{cor:favorable-corridor}
Let \(\rho(s)\to\rho^+\) as \(s\to+\infty\).  If
\[
  (c+\varepsilon)n\leq x_n
  \leq(c_{\rm front}-\varepsilon)n
\]
for every \(n\), where \(\varepsilon>0\), then
\[
  \rho(x_n-cn)\longrightarrow\rho^+.
\]
The corridor is nonempty whenever \(c+2\varepsilon\leq c_{\rm front}\).
\end{corollary}

\begin{proof}
The first conclusion follows by composing the far-right limit of \(\rho\)
with Lemma~\ref{lem:moving-coordinate}(iii).  The nonemptiness condition is
the inequality
\[
  c+\varepsilon\leq c_{\rm front}-\varepsilon.
\]
\end{proof}

Thus convergence to a far-right favorable equilibrium belongs in a
one-sided corridor
\[
  (c+\varepsilon)n
  \leq x
  \leq(c_{\rm front}-\varepsilon)n.
\]
A symmetric band \(|x-cn|\leq\eta n\) contains
\(x=(c-\eta)n\), whose moving coordinate tends to \(-\infty\).

\section{Corrected global research questions}
\label{sec:global-conjectures}

This section records precise global statements suggested by the preceding
estimates.  They are conjectures, not consequences of the local and
upper-barrier arguments alone.

Let
\[
  R_i:=1+\rho_i^+,
  \qquad
  \M_i(\mu):=\int_{\R}k_i(z)e^{\mu z}\,dz,
\]
and define the empty-habitat rightward speeds
\begin{equation}\label{eq:empty-speeds}
  c_i^*
  :=
  \inf_{\mu>0}
  \frac{\log(R_i\M_i(\mu))}{\mu}.
\end{equation}
For weak competition, define the favorable rare-species multipliers
\[
  \widetilde R_1
  :=
  \frac{R_1}{1+\rho_1^+\alpha_1},
  \qquad
  \widetilde R_2
  :=
  \frac{R_2}{1+\rho_2^+\alpha_2},
\]
the associated candidate invasion speeds
\begin{equation}\label{eq:reduced-speeds}
  \widetilde c_i^*
  :=
  \inf_{\mu>0}
  \frac{\log(\widetilde R_i\M_i(\mu))}{\mu},
  \qquad
  c_{\rm coex}:=\min_i\widetilde c_i^*.
\end{equation}
Identifying the nonlinear spreading speed with one of these variational
quantities requires the usual irreducibility and linear-determinacy
hypotheses.

\begin{conjecture}[Fast-habitat extinction]
\label{conj:global-extinction}
Assume the hypotheses of Section~2, the required exponential moments, and
compactly supported initial data in \([0,1]^2\).  If
\[
  c>\max\{c_1^*,c_2^*\},
\]
then
\[
  \sup_{x\in\R}\max\{u_n(x),v_n(x)\}\longrightarrow0.
\]
\end{conjecture}

\begin{conjecture}[Intermediate-speed exclusion]
\label{conj:global-exclusion}
Assume \(c_1^*<c<c_2^*\) and \(v_0\not\equiv0\).  Under suitable scalar
persistence and linear-determinacy hypotheses,
\[
  \sup_{x\in\R}u_n(x)\longrightarrow0,
\]
and, for every
\[
  0<\varepsilon<\frac{c_2^*-c}{2},
\]
\[
  \sup_{(c+\varepsilon)n\leq x\leq(c_2^*-\varepsilon)n}
  |v_n(x)-1|\longrightarrow0.
\]
A complete spreading statement must specify the rightward front and its
behavior ahead of \(c_2^*\); if kernels are asymmetric, the leftward and
rightward speeds must be distinguished.
\end{conjecture}

\begin{conjecture}[Weak-competition coexistence]
\label{conj:global-coexistence}
Assume \(0<\alpha_1,\alpha_2<1\), both initial components are nonzero, and
the required irreducibility and linear-determinacy hypotheses hold.  If
\[
  0<c<c_{\rm coex},
\]
then, for every
\[
  0<\varepsilon<\frac{c_{\rm coex}-c}{2},
\]
\[
  \sup_{(c+\varepsilon)n\leq x\leq(c_{\rm coex}-\varepsilon)n}
  \max\{|u_n(x)-u^*|,\ |v_n(x)-v^*|\}
  \longrightarrow0.
\]
Equivalently, the corresponding fixed-speed observer statement should
quantify only over
\[
  c<s<c_{\rm coex}.
\]
\end{conjecture}

\begin{remark}
The reduced quantities in \eqref{eq:reduced-speeds} are natural sufficient
thresholds because they arise from invasion of the two resident-only states.
They should not be presented as exact nonlinear speeds without a separate
linear-determinacy theorem.  In particular, the empty-habitat speeds
\eqref{eq:empty-speeds} alone do not encode the competition coefficients.
\end{remark}

\section{A reusable generalization}

The same local argument applies to the more general response
\begin{equation}\label{eq:general-response}
  G_i(s;u,v)
  :=
  \frac{R_i(s)u}
       {1+b_i(s)(u+\alpha_i v)},
  \qquad
  R_i(s)>0,\quad b_i(s)\geq0.
\end{equation}
Positivity and competitive order follow immediately.  A sufficient
condition for \(G_i(s;u,v)\leq1\) whenever \(0\leq u\leq1\) is
\[
  R_i(s)\leq1+b_i(s).
\]
Indeed,
\[
  1+b_i(s)(u+\alpha_i v)-R_i(s)u
  =
  1-\bigl(R_i(s)-b_i(s)\bigr)u+b_i(s)\alpha_i v
\]
If \(R_i(s)\leq b_i(s)\), the right-hand side is at least one.  Otherwise
\(0<R_i(s)-b_i(s)\leq1\), so
\(\bigl(R_i(s)-b_i(s)\bigr)u\leq1\), and the right-hand side is again
nonnegative.
The model \eqref{eq:corrected-response} is the special case
\[
  R_i(s)=1+\rho_i(s),
  \qquad
  b_i(s)=[\rho_i(s)]_+.
\]
This formulation separates low-density growth, density dependence, and
interspecific competition.  It also permits different environments for the
two species, asymmetric kernels, directional moments in higher dimensions,
and habitat displacements that vary with time.

\appendix

\section{Audit of the original proposal}
\label{app:audit}

This appendix concerns the CMBE 2026 proposal
\emph{Spatial Dynamics of an Integro-Difference Lotka--Volterra Model in a
Shifting Habitat}.  The purpose is diagnostic: for each defect, we give the
shortest mathematical reason that the original statement cannot stand and
then record a repair.

\subsection{The printed response is not positive on the stated state space}

The response printed in the proposal is
\begin{equation}\label{eq:printed-response}
  F(\rho,\alpha;u,v)
  :=
  \frac{(1+\rho)u}{1+\rho(u+\alpha v)}.
\end{equation}
The hypotheses of Theorems 2.1--2.3 allow arbitrary bounded nonnegative
initial data.  They do not impose \(u+\alpha v\leq1\).  Under the permitted
values
\[
  \rho=-\frac12,\qquad \alpha=0,\qquad u=3,\qquad v=0,
\]
one obtains
\[
  F\!\left(-\frac12,0;3,0\right)
  =
  \frac{(1/2)3}{1-(1/2)3}
  =-3.
\]
Thus the local update can send an admissible nonnegative state to a negative
value.  The model, as stated, is not a positive population dynamical system.

\paragraph{Repair.}
Either add and prove an invariant state-space restriction strong enough to
keep every denominator positive, or change the response.  The correction
\eqref{eq:corrected-response} gives positivity on the entire positive cone
and is therefore the cleaner repair.

\subsection{The simplex remark does not restore competition}

The proposal contains a red editorial sentence suggesting
\[
  1-u-\alpha v\geq0.
\]
This is not included as a theorem hypothesis, and no invariance proof is
given.  More importantly, even if the inequality is imposed, the direction
of competition remains wrong in unfavorable habitat.  For
\[
  \rho=-\frac12,\qquad \alpha=1,\qquad u=\frac12,
\]
both \(v=0\) and \(v=\frac12\) satisfy \(u+\alpha v\leq1\), but
\[
  F\!\left(-\frac12,1;\frac12,0\right)=\frac13
  <
  \frac12
  =
  F\!\left(-\frac12,1;\frac12,\frac12\right).
\]
Increasing the competitor increases the focal-species response.  In fact,
where the denominator is positive,
\[
  \frac{\partial F}{\partial v}
  =
  -\frac{(1+\rho)u\rho\alpha}
         {(1+\rho(u+\alpha v))^2},
\]
which is positive when \(-1<\rho<0\), \(u>0\), and \(\alpha>0\).

\paragraph{Repair.}
Density dependence must remain nonnegative.  Replacing \(\rho\) by
\([\rho]_+\) in the denominator yields
\eqref{eq:corrected-response}; the unfavorable response then becomes
subcritical linear growth.

\subsection{Theorem 2.1 lacks a valid invariant state space}

Theorem 2.1 assumes only that the initial pair is nonnegative, bounded, and
compactly supported.  The calculation above uses data allowed by exactly
those assumptions.  Consequently, the theorem is not formulated on a state
space preserved by its own update, and the comparison principles invoked in
the proposal are unavailable.

This defect is logically prior to the speed calculation.  Even a correct
formula for \(c_i^*\) cannot repair a nonlinear map that leaves the positive
cone.

\paragraph{Repair.}
Use the corrected response and initial values in the invariant unit square
\([0,1]^2\), as in Conjecture~\ref{conj:global-extinction}.  A proof must
then combine the upper exponential estimate with decay in the unfavorable
region; the barrier alone proves only the upper-front part.

\subsection{Theorem 2.2 has a missing-species counterexample}

Theorem 2.2 assumes that the pair \((u_0,v_0)\) has nonempty compact support,
but it does not assume \(v_0\not\equiv0\).  Choose
\[
  u_0\not\equiv0,\qquad v_0\equiv0.
\]
Then the pair satisfies the printed initial-data hypothesis.  Since the
second response has a factor \(v_n(y)\), induction gives
\[
  v_n\equiv0\qquad\text{for all }n.
\]
Therefore
\[
  \liminf_{n\to\infty}
  \inf_{|x-cn|\leq\eta n}v_n(x)=0
\]
for every \(\eta>0\), contradicting the asserted strict positivity.

\paragraph{Repair.}
Any persistence conclusion for species two must assume
\(v_0\not\equiv0\), together with the positivity or irreducibility
conditions needed to spread nonzero mass.

\subsection{Theorem 2.2 uses the wrong persistence region}

The symmetric region in the original theorem is
\[
  |x-cn|\leq\eta n.
\]
Its left boundary is \(x=(c-\eta)n\), so in habitat coordinates
\[
  x-cn=-\eta n\longrightarrow-\infty.
\]
The left edge therefore samples the unfavorable limit \(\rho^-\), while the
right edge samples the favorable limit.  A uniform lower bound over the
entire symmetric band cannot be justified by convergence to the favorable
single-species equilibrium.

\paragraph{Repair.}
Use a one-sided interior corridor
\[
  (c+\varepsilon)n
  \leq x
  \leq(c_2^*-\varepsilon)n.
\]
Every sequence in this corridor has habitat coordinate tending to
\(+\infty\).  The phrase ``spreads asymptotically with speed \(c_2^*\)''
must also be replaced by a quantified rightward spreading statement.  Since
the proposal does not assume symmetric kernels and separately defines
leftward moments, leftward and rightward speeds cannot be conflated.

\subsection{Theorem 2.3 quantifies over observers behind the habitat}

Theorem 2.3 asserts convergence to the positive favorable equilibrium for
every \(0<c'<c^\dagger\), where \(c^\dagger>c\).  This interval includes all
\(c'\) with \(0<c'<c\).  For such an observer,
\[
  c'n-cn=(c'-c)n\longrightarrow-\infty.
\]
Hence the observer moves linearly far into the unfavorable side of the
habitat.  The claimed convergence to the far-right positive coexistence
equilibrium uses the wrong quantifier range.

\paragraph{Repair.}
The observer speed must satisfy
\[
  c<c'<c_{\rm coex},
\]
or the conclusion should be stated uniformly in the corresponding
one-sided corridor.

\subsection{Theorem 2.3 uses only empty-habitat speeds}

The threshold
\[
  c<\min\{c_1^*,c_2^*\}
\]
uses growth at the empty state.  Coexistence, however, requires each species
to invade a state occupied by the other species.  At resident density one,
the rare-species multiplier is
\[
  \widetilde R_i
  =
  \frac{1+\rho_i^+}{1+\rho_i^+\alpha_i},
\]
not \(1+\rho_i^+\).  For every
\(0<\alpha_i<1\), this multiplier is greater than one but strictly smaller
than the empty-state multiplier.  The associated invasion front can
therefore be slower than the empty-habitat front.

\paragraph{Repair.}
Formulate a conservative coexistence corridor using the reduced speeds in
\eqref{eq:reduced-speeds}, and state separately the irreducibility and
linear-determinacy conditions under which those variational speeds are
exact.

\subsection{The strong-competition conclusion does not follow from speed
ordering}

The summary states that under strong competition the faster species excludes
the other, even when both persist alone.  Suppose, in the favorable
homogeneous environment, that
\[
  \alpha_1>1,\qquad \alpha_2>1.
\]
By Proposition~\ref{prop:invasion-multiplier}, each rare-species multiplier
at the other resident-only equilibrium is less than one:
\[
  \frac{1+\rho^+}{1+\rho^+\alpha_i}<1.
\]
Thus both semi-trivial equilibria resist invasion.  This is the local
signature of bistability.  Ordering the two empty-habitat speeds does not
select a unique winner independently of initial data.

\paragraph{Repair.}
A strong-competition theorem must specify invasion eigenvalues, wave
selection, and initial-data conditions.  It may yield conditional exclusion
or a bistable traveling front, but speed ordering alone is insufficient.

\subsection{The PINN convergence sentence is not an analytical consequence}

The proposal suggests that training loss tending to zero should imply
convergence to the analytical traveling wave and consistent recovery of its
speed.  A small sampled residual does not by itself imply either conclusion.
One must specify at least:
\begin{enumerate}[label=(\roman*)]
  \item the function space and norm in which the residual is controlled;
  \item compactness or coercivity for the network sequence;
  \item consistency of quadrature for the nonlocal integral;
  \item stability and uniqueness of the target wave after fixing translation;
  \item identifiability of the wave speed.
\end{enumerate}
Without these ingredients, the PINN is a numerical proposal rather than a
convergence theorem.

\paragraph{Repair.}
State the neural-network component as a computational method and validate it
numerically.  A mathematical convergence claim should be deferred until a
separate approximation theorem supplies the missing hypotheses.

\subsection{The status and notation are internally inconsistent}

Section 2 says that the three theorems are results ``we propose to
establish,'' whereas the summary says ``we prove'' and the conclusion says
``we established.''  These descriptions cannot all be correct in the same
document.  In addition, both species equations use \(\rho_1\).  This is
legitimate only if a common environmental response is intended; otherwise
the second occurrence should be \(\rho_2\).

\paragraph{Repair.}
Until the global arguments are complete, use ``we conjecture'' or ``we aim
to prove.''  State explicitly whether the two species share one
environmental response or have species-specific functions.

\section{Summary of repairs}

\begin{center}
\begin{tabularx}{0.92\textwidth}{@{}>{\raggedright\arraybackslash}p{0.27\textwidth}Y@{}}
\toprule
Original issue & Required repair\\
\midrule
Signed density dependence &
Use nonnegative density dependence, for example
\([\rho]_+\), or prove a genuinely invariant alternative state space.\\
Bounded nonnegative data &
Work in the invariant unit square and require every persisting species to be
initially nonzero.\\
Symmetric persistence band &
Use a one-sided corridor strictly ahead of the habitat.\\
Observer range \(0<c'<c^\dagger\) &
Replace it by \(c<c'<c_{\rm coex}\).\\
Empty-habitat coexistence threshold &
Use invasion multipliers and prove linear determinacy at the semi-trivial
states.\\
Strong competition &
Add invasion and initial-data hypotheses; allow bistability.\\
PINN convergence &
Treat it as numerical methodology unless a separate approximation theorem is
proved.\\
\bottomrule
\end{tabularx}
\end{center}

\begin{thebibliography}{9}

\bibitem{LewisMarculisShen2018}
M.~A. Lewis, N.~G. Marculis, and Z. Shen,
\href{https://www.math.ualberta.ca/~mlewis/Publications%202018/Lewis_Marculis_Shen.pdf}
{Integrodifference equations in the presence of climate change:
persistence criterion, travelling waves and inside dynamics},
\emph{Journal of Mathematical Biology} 77 (2018), 1649--1687.

\bibitem{WuWangZou2019}
C. Wu, Y. Wang, and X. Zou,
\href{https://www.uwo.ca/apmaths/faculty/zou/repr/jde19-2.pdf}
{Spatial-temporal dynamics of a Lotka--Volterra competition model with
nonlocal dispersal under shifting environment},
\emph{Journal of Differential Equations} 267 (2019), 4890--4921.

\bibitem{KongRawalShen2015}
L. Kong, N. Rawal, and W. Shen,
\href{https://arxiv.org/abs/1410.0317}
{Spreading speeds and linear determinacy for two species competition
systems with nonlocal dispersal in periodic habitats},
\emph{Mathematical Modelling of Natural Phenomena} 10 (2015), 113--141.

\end{thebibliography}

\end{document}
